\chapter{Directed Graphs}
%%%%%%%%%%%%%%%%%%%%%%%%%%%%%%%%%%%%%%%%%%%%%%%%%%%%%%%%%%%%%%%%%%%%%%%%%%%%%%%
\section{Directed Graphs}

% Generated by scripts/blueprint_restate.py from TCSlib.GraphTheory.Core.DirectedGraphs.
% Each body is the STATEMENT only, taken from the declaration's Lean docstring:
% the one-line summary plus the verbatim book statement.  Proof, proof plan,
% significance commentary and formalisation notes are excluded by construction.

\begin{definition}[Edge connectivity $\kappa'(G)$]
\lean{edgeConnectivity}
\label{edgeConnectivity}
\leanok
\uses{SimpleGraph.IsEdgeCut}
Edge connectivity \texttt{$\kappa$'(G)}.

\emph{We \dots{} define the edge connectivity
\texttt{$\kappa$'(G)} of \texttt{G} to be the minimum \texttt{k} for which \texttt{G} has a \texttt{k}-edge cut.  If \texttt{G} is
trivial, \texttt{$\kappa$'(G)} is defined to be zero. \dots{} \texttt{G} is said to be
\texttt{k}-edge-connected if
\texttt{$\kappa$'(G) $\ge$ k}.}
\end{definition}

\begin{definition}[The bridge \texttt{Digraph $\to$ Quiver}: \texttt{Quiver.\{0\} V} has \texttt{Hom : V $\to$ V $\to$ Prop},\dots]
\lean{Digraph.toQuiver}
\label{Digraph.toQuiver}
\leanok
The bridge \texttt{Digraph $\to$ Quiver}: \texttt{Quiver.\{0\} V} has \texttt{Hom : V
$\to$ V $\to$ Prop}, exactly \texttt{Digraph.Adj}.
\end{definition}

\begin{definition}[A directed path: a \texttt{Quiver.Path} with no repeated vertex]
\lean{Digraph.IsDirectedPath}
\label{Digraph.IsDirectedPath}
\leanok
\uses{Digraph.toQuiver}
A directed \textbf{path}: a \texttt{Quiver.Path} with no repeated vertex.

\emph{A directed walk in \texttt{D} is a finite
non-null sequence \texttt{W = (v$_{0}$, a$_{1}$, v$_{1}$, \dots{}, a\_k, v\_k)}, whose
terms are alternately
vertices and arcs, such that, for \texttt{i = 1, 2, \dots{}, k}, the arc
\texttt{a$_{i}$} has head \texttt{v$_{i}$} and
tail \texttt{v\_\{i-1\}}. \dots{} A directed trail is a directed walk that is a trail;
directed
paths, directed cycles and directed tours are similarly defined.}
\end{definition}

\begin{definition}[A directed cycle: a positive-length closed walk whose vertices (bar the repeated endpoint) are distinct]
\lean{Digraph.IsDirectedCycle}
\label{Digraph.IsDirectedCycle}
\leanok
\uses{Digraph.toQuiver}
A directed \textbf{cycle}: a positive-length closed walk whose vertices (bar the
repeated endpoint)
are distinct.

The book defines it only by analogy --- \emph{directed
paths, directed cycles and directed tours are similarly defined} --- so the unfolding
here (a closed directed walk of positive length whose vertices, apart from the
coinciding endpoints, are distinct) follows \S{}1.4's cycle.
\end{definition}

\begin{definition}[$v$ is reachable from $u$]
\lean{Digraph.Reachable}
\label{Digraph.Reachable}
\leanok
\texttt{v} is \textbf{reachable} from \texttt{u}.

\emph{If there is a directed \texttt{(u, v)}-path in
\texttt{D}, vertex \texttt{v} is said to be reachable from vertex \texttt{u} in \texttt{D}.}
\end{definition}

\begin{definition}[$D$ is diconnected: every vertex reaches every other]
\lean{Digraph.Diconnected}
\label{Digraph.Diconnected}
\leanok
\uses{Digraph.Reachable}
\texttt{D} is \textbf{diconnected}: every vertex reaches every other.

\emph{Two vertices are diconnected in \texttt{D} if
each is reachable from the other.  As in the case of connection in graphs,
diconnection is an equivalence relation on the vertex set of \texttt{D}. The subdigraphs
\texttt{D[V$_{1}$], D[V$_{2}$], \dots{}, D[V\_m]} induced by the resulting partition \texttt{(V$_{1}$, V$_{2}$, \dots{}, V\_m)} of
\texttt{V(D)} are called the dicomponents of \texttt{D}.  A digraph \texttt{D} is diconnected if it has
exactly one dicomponent.}
\end{definition}

\begin{definition}[$d^{-}(v)$, the indegree]
\lean{Digraph.indegree}
\label{Digraph.indegree}
\leanok
\texttt{d$^{-}$(v)}, the \textbf{indegree}.

\emph{The indegree \texttt{d$^{-}$\_D(v)} of a vertex \texttt{v}
in \texttt{D} is the number of arcs with head \texttt{v}.}
\end{definition}

\begin{definition}[$d^{+}(v)$, the outdegree]
\lean{Digraph.outdegree}
\label{Digraph.outdegree}
\leanok
\texttt{d$^{+}$(v)}, the \textbf{outdegree}.

\emph{\dots{}the outdegree \texttt{d$^{+}$\_D(v)} of \texttt{v} is the
number of arcs with tail \texttt{v}.}
\end{definition}

\begin{definition}[$\varepsilon$, the number of arcs]
\lean{Digraph.arcCount}
\label{Digraph.arcCount}
\leanok
\texttt{$\varepsilon$}, the number of \textbf{arcs}.

\emph{Throughout this chapter, \texttt{D} will
denote a digraph and \texttt{G} its underlying graph.  This is a useful convention; it
allows us, for example, to denote the vertex set of \texttt{D} by \texttt{V} (since
\texttt{V = V(G)}), and the numbers of vertices and arcs in \texttt{D} by \texttt{$\nu$} and \texttt{$\varepsilon$},
respectively.}
\end{definition}

\begin{definition}[$\delta^{-}$, the minimum indegree]
\lean{Digraph.minIndegree}
\label{Digraph.minIndegree}
\leanok
\uses{Digraph.indegree}
\texttt{$\delta$$^{-}$}, the minimum indegree.

\emph{We denote the minimum and maximum
indegrees and outdegrees in \texttt{D} by \texttt{$\delta$$^{-}$(D)},
\texttt{$\Delta$$^{-}$(D)}, \texttt{$\delta$$^{+}$(D)} and \texttt{$\Delta$$^{+}$(D)},
respectively.}
\end{definition}

\begin{definition}[$\delta^{+}$, the minimum outdegree]
\lean{Digraph.minOutdegree}
\label{Digraph.minOutdegree}
\leanok
\uses{Digraph.outdegree}
\texttt{$\delta$$^{+}$}, the minimum outdegree.

The same sentence as for \texttt{$\delta$$^{-}$}: \emph{We denote the
minimum and maximum indegrees and outdegrees in \texttt{D} by
\texttt{$\delta$$^{-}$(D)}, \texttt{$\Delta$$^{-}$(D)}, \texttt{$\delta$$^{+}$(D)}
and \texttt{$\Delta$$^{+}$(D)}, respectively.}
\end{definition}

\begin{definition}[The converse $\breve{D}$: reverse every arc]
\lean{Digraph.converse}
\label{Digraph.converse}
\leanok
The \textbf{converse} \texttt{$\breve{D}$}: reverse every arc.

\emph{The converse \texttt{$\breve{D}$} of \texttt{D} is
the digraph obtained from \texttt{D} by reversing the orientation of each arc.}
\end{definition}

\begin{definition}[$D$ is strict: loopless (\texttt{Digraph} already forbids parallel same-direction arcs)]
\lean{Digraph.IsStrict}
\label{Digraph.IsStrict}
\leanok
\texttt{D} is \textbf{strict}: loopless (\texttt{Digraph} already forbids parallel same-direction arcs).

\emph{A digraph is strict if it has no loops
and no two arcs with the same ends have the same orientation.}
\end{definition}

\begin{definition}[$D$ is an orientation of $G$ --- the chapter's keystone gap]
\lean{Digraph.IsOrientationOf}
\label{Digraph.IsOrientationOf}
\leanok
\texttt{D} is an \textbf{orientation} of \texttt{G} --- the chapter's keystone gap.

\emph{With each digraph \texttt{D} we can associate
a graph \texttt{G} on the same vertex set; corresponding to each arc of \texttt{D} there
is an
edge of \texttt{G} with the same ends. This graph is the underlying graph of \texttt{D}.
Conversely, given any graph \texttt{G}, we can obtain a digraph from \texttt{G} by
specifying,
for each link, an order on its ends.  Such a digraph is called an orientation of
\texttt{G}.}
\end{definition}

\begin{definition}[$D$ is a tournament: an orientation of the complete graph]
\lean{Digraph.IsTournament}
\label{Digraph.IsTournament}
\leanok
\texttt{D} is a \textbf{tournament}: an orientation of the complete graph.

\emph{An orientation of a complete graph is
called a tournament.  The tournaments on four vertices are shown in figure 10.4.
Each can be regarded as indicating the results of games in a round-robin tournament
between four players; for example, the first tournament in figure 10.4 shows that
one player has won all three games and that the other three have each won one.}
\end{definition}

\begin{definition}[B\&M's $(S, T)$: arcs with tail in $S$, head in $T$]
\lean{Digraph.arcsBetween}
\label{Digraph.arcsBetween}
\leanok
B\&M's \texttt{(S, T)}: arcs with tail in \texttt{S}, head in \texttt{T}.

\emph{If \texttt{S} and \texttt{T} are subsets of \texttt{V}, we
denote by \texttt{(S, T)} the set of arcs of \texttt{D} that have their tails in
\texttt{S} and their
heads in \texttt{T}.}
\end{definition}

\begin{definition}[$D$ is $k$-arc-connected: every nonempty proper cut has $\ge k$ outgoing arcs]
\lean{Digraph.IsKArcConnected}
\label{Digraph.IsKArcConnected}
\leanok
\uses{Digraph.arcsBetween}
\texttt{D} is \textbf{\texttt{k}-arc-connected}: every nonempty proper cut has \texttt{$\ge$ k} outgoing arcs.

\emph{A nontrivial digraph \texttt{D} is
\texttt{k}-arc-connected if, for every nonempty proper subset \texttt{S} of \texttt{V},
\texttt{|(S, $\bar{S}$)| $\ge$ k}.}
\end{definition}

\begin{definition}[The induced subdigraph on $S$ (kept on the same carrier to avoid subtype juggling)]
\lean{Digraph.induce}
\label{Digraph.induce}
\leanok
The \textbf{induced subdigraph} on \texttt{S} (kept on the same carrier to avoid subtype
juggling).

\emph{A digraph \texttt{D'} is a subdigraph of \texttt{D}
if \texttt{V(D') $\subseteq$ V(D)}, \texttt{A(D') $\subseteq$ A(D)} and
\texttt{$\psi$\_\{D'\}} is the restriction of \texttt{$\psi$\_D} to
\texttt{A(D')}.  The terminology and notation for subdigraphs is similar to that used for
subgraphs.}
\end{definition}

\begin{definition}[The list of arcs traversed by a directed walk]
\lean{Digraph.arcsOf}
\label{Digraph.arcsOf}
\uses{Digraph.toQuiver}
The list of arcs traversed by a directed walk.
\end{definition}

\begin{definition}[A directed trail: no repeated arc]
\lean{Digraph.IsDirectedTrail}
\label{Digraph.IsDirectedTrail}
\uses{Digraph.arcsOf, Digraph.toQuiver}
A directed \textbf{trail}: no repeated arc.

\emph{A directed trail is a directed walk
that is a trail} --- i.e. one whose arcs are all distinct.
\end{definition}

\begin{definition}[A directed Euler tour: a closed directed trail using every arc]
\lean{Digraph.IsDirectedEulerTour}
\label{Digraph.IsDirectedEulerTour}
\leanok
\uses{Digraph.IsDirectedTrail, Digraph.arcsOf, Digraph.toQuiver}
A directed \textbf{Euler tour}: a closed directed trail using every arc.

\emph{A directed Euler tour of \texttt{D}
is a directed tour that traverses each arc of \texttt{D} exactly once.}
\end{definition}

\begin{definition}[$D$ is unilateral: any two vertices are comparable by reachability]
\lean{Digraph.IsUnilateral}
\label{Digraph.IsUnilateral}
\leanok
\uses{Digraph.Reachable}
\texttt{D} is \textbf{unilateral}: any two vertices are comparable by reachability.

\emph{A digraph \texttt{D} is unilateral
if, for any two vertices \texttt{u} and \texttt{v}, either \texttt{v} is reachable from
\texttt{u} or \texttt{u} is
reachable from \texttt{v}.}
\end{definition}

\begin{definition}[The 0/1 adjacency matrix of $D$]
\lean{Digraph.adjMatrix}
\label{Digraph.adjMatrix}
\leanok
The 0/1 \textbf{adjacency matrix} of \texttt{D}.

\emph{Let \texttt{v$_{1}$, v$_{2}$, \dots{}, v\_$\nu$} be the
vertices of a digraph \texttt{D}. The adjacency matrix of \texttt{D} is the
\texttt{$\nu$ $\times$ $\nu$} matrix
\texttt{A = [a\_ij]} in which \texttt{a\_ij} is the number of arcs of \texttt{D} with tail \texttt{v$_{i}$} and head
\texttt{v$_{j}$}.}
\end{definition}

\begin{definition}[The condensation $\hat{D}$: dicomponents contracted]
\lean{Digraph.condensation}
\label{Digraph.condensation}
\leanok
\uses{Digraph.toQuiver}
The \textbf{condensation} \texttt{$\hat{D}$}: dicomponents contracted.

\emph{Let \texttt{D$_{1}$, D$_{2}$, \dots{}, D\_m} be the
dicomponents of \texttt{D}. The condensation \texttt{$\hat{D}$} of \texttt{D} is a
directed graph with \texttt{m}
vertices \texttt{w$_{1}$, w$_{2}$, \dots{}, w\_m}; there is an arc in \texttt{$\hat{D}$}
with tail \texttt{w$_{i}$} and head \texttt{w$_{j}$} if
and only if there is an arc in \texttt{D} with tail in \texttt{D$_{i}$} and head in
\texttt{D$_{j}$}.}
\end{definition}

\begin{definition}[dirDist]
\lean{Digraph.dirDist}
\label{Digraph.dirDist}
\leanok
\uses{Digraph.toQuiver}
\textbf{Directed distance} \texttt{d(u,v)}.

\emph{In a diconnected digraph \texttt{D}, the
length of a shortest directed \texttt{(u, v)}-path is denoted by \texttt{d\_D(u, v)} and
is called
the distance from \texttt{u} to \texttt{v}; the directed diameter of \texttt{D} is the
maximum distance
from any one vertex of \texttt{D} to any other.}
\end{definition}

\begin{definition}[dirDiameter]
\lean{Digraph.dirDiameter}
\label{Digraph.dirDiameter}
\leanok
\uses{Digraph.dirDist}
\textbf{Directed diameter}.

\emph{\dots{}the directed diameter of \texttt{D} is the
maximum distance from any one vertex of \texttt{D} to any other.}
\end{definition}

\begin{definition}[The associated digraph $D(G)$: each edge becomes two opposite arcs. \texttt{G.Adj} is already symmetric,\dots]
\lean{SimpleGraph.associatedDigraph}
\label{SimpleGraph.associatedDigraph}
\leanok
The \textbf{associated digraph} \texttt{D(G)}: each edge becomes two opposite arcs.
\texttt{G.Adj} is already symmetric, so it \emph{is} the arc relation.

\emph{The associated digraph \texttt{D(G)}
of a graph \texttt{G} is the digraph obtained when each edge \texttt{e} of \texttt{G} is
replaced by two
oppositely oriented arcs with the same ends as \texttt{e}.}
\end{definition}

\begin{definition}[A primitive $\mathbb{N}$-matrix: some power is entrywise positive]
\lean{Matrix.IsPrimitive}
\label{Matrix.IsPrimitive}
\leanok
A \textbf{primitive} $\mathbb{N}$-matrix: some power is entrywise positive.

\emph{A real matrix \texttt{R} is called primitive
if \texttt{R$^{k}$ > 0} for some \texttt{k}.}
\end{definition}

\begin{definition}[The de Bruijn digraph $D_n$: vertices are $(n-1)$-bit strings, arcs are left-shifts]
\lean{deBruijnDigraph}
\label{deBruijnDigraph}
\leanok
The \textbf{de Bruijn digraph} \texttt{D\_n}: vertices are \texttt{(n-1)}-bit strings,
arcs are left-shifts.

\emph{We define a digraph \texttt{D\_n} as follows:
the vertices of \texttt{D\_n} are the \texttt{(n-1)}-digit binary numbers
\texttt{p$_{1}$ p$_{2}$ \dots{} p\_\{n-1\}} with
\texttt{p$_{i}$ = 0} or \texttt{1}.  There is an arc with tail \texttt{p$_{1}$ p$_{2}$ \dots{} p\_\{n-1\}} and head
\texttt{q$_{1}$ q$_{2}$ \dots{} q\_\{n-1\}} if and only if \texttt{p\_\{i+1\} = q$_{i}$} for \texttt{1 $\le$ i $\le$ n-2}; in other words,
all arcs are of the form \texttt{(p$_{1}$ p$_{2}$ \dots{} p\_\{n-1\}, p$_{2}$ p$_{3}$
\dots{} p\_n)}. In addition, each arc
\texttt{(p$_{1}$ p$_{2}$ \dots{} p\_\{n-1\}, p$_{2}$ p$_{3}$ \dots{} p\_n)} of \texttt{D\_n} is assigned the label \texttt{p$_{1}$ p$_{2}$ \dots{} p\_n}.}
\end{definition}

\begin{theorem}[Counting the orientations of a finite simple graph]
\lean{card_orientations}
\label{card_orientations}
\leanok
\uses{Digraph.IsOrientationOf}
Let $G$ be a simple graph on a finite vertex set. Then the number of orientations of $G$
is $2^{\abs{E(G)}}$, where $\abs{E(G)}$ denotes the number of edges of $G$; that is,
each edge may be directed in either of its two ways, and every such choice yields a
distinct orientation.
\end{theorem}

\begin{theorem}[Degree sums count the arcs of a digraph]
\lean{sum_indegree_eq_arcCount}
\label{sum_indegree_eq_arcCount}
\leanok
\uses{Digraph.arcCount, Digraph.indegree, Digraph.outdegree}
Let $D$ be a digraph on a finite vertex set $V$, and write $\varepsilon$ for its number
of arcs. Then both the sum of the indegrees and the sum of the outdegrees over all
vertices equal $\varepsilon$:
\[
\sum_{v \in V} d^{-}(v) \;=\; \varepsilon \;=\; \sum_{v \in V} d^{+}(v).
\]
\end{theorem}

\begin{theorem}[Acyclic digraphs have a source]
\lean{exists_indegree_zero_of_acyclic}
\label{exists_indegree_zero_of_acyclic}
\leanok
\uses{Digraph.IsDirectedCycle, Digraph.indegree, Digraph.toQuiver}
Let $D$ be a digraph on a finite, nonempty vertex set that contains no directed cycle.
Then $D$ has a vertex $v$ whose indegree vanishes, $d^{-}(v) = 0$; that is, no arc of
$D$ has $v$ as its head.
\end{theorem}

\begin{theorem}[Topological ordering of an acyclic digraph]
\lean{exists_topological_ordering}
\label{exists_topological_ordering}
\leanok
\uses{Digraph.IsDirectedCycle, Digraph.toQuiver}
Let $D$ be a digraph on a finite vertex set $V$ with $n = |V|$, and suppose $D$ is
acyclic: it contains no directed cycle, that is, no closed directed walk of positive
length whose vertices, apart from the coinciding endpoints, are distinct. Then $V$
admits a topological ordering: there is a bijection $f \colon V \to \{0, 1, \dots,
n-1\}$ such that $f(u) < f(v)$ for every arc of $D$ from $u$ to $v$.
\end{theorem}

\begin{theorem}[The converse is an involution]
\lean{converse_converse}
\label{converse_converse}
\leanok
\uses{Digraph.converse}
Let $D$ be a digraph on a vertex set $V$. Reversing every arc twice recovers the
original digraph, so the converse of the converse of $D$ equals $D$; that is,
$\breve{\breve{D}} = D$.
\end{theorem}

\begin{theorem}[Degrees swap under the converse digraph]
\lean{converse_indegree}
\label{converse_indegree}
\leanok
\uses{Digraph.converse, Digraph.indegree, Digraph.outdegree}
Let $D$ be a digraph on a finite vertex set, and let $\breve{D}$ be its converse,
obtained by reversing the orientation of every arc. Then for each vertex $v$ the
indegree and outdegree interchange:
\[
d^{-}_{\breve{D}}(v) = d^{+}_{D}(v)
\qquad\text{and}\qquad
d^{+}_{\breve{D}}(v) = d^{-}_{D}(v),
\]
where $d^{-}$ counts the arcs with head $v$ and $d^{+}$ the arcs with tail $v$.
\end{theorem}

\begin{theorem}[Reachability in the converse digraph]
\lean{converse_reachable}
\label{converse_reachable}
\leanok
\uses{Digraph.Reachable, Digraph.converse}
Let $D$ be a digraph with vertices $u$ and $v$, and let $\breve{D}$ denote its converse,
the digraph obtained by reversing the orientation of every arc. Then $v$ is reachable
from $u$ in $\breve{D}$ if and only if $u$ is reachable from $v$ in $D$.
\end{theorem}

\begin{theorem}[Every finite acyclic digraph has a sink]
\lean{exists_outdegree_zero_of_acyclic}
\label{exists_outdegree_zero_of_acyclic}
\leanok
\uses{Digraph.IsDirectedCycle, Digraph.outdegree, Digraph.toQuiver}
Let $D$ be a digraph on a finite, nonempty vertex set $V$. If $D$ contains no directed
cycle, then there is a vertex $v \in V$ whose outdegree vanishes, $d^{+}(v) = 0$.
\end{theorem}

\begin{theorem}[A directed path as long as the minimum degree]
\lean{exists_directedPath_length_ge_maxMinDegree}
\label{exists_directedPath_length_ge_maxMinDegree}
\leanok
\uses{Digraph.IsDirectedPath, Digraph.IsStrict, Digraph.minIndegree, Digraph.minOutdegree, Digraph.toQuiver}
Let $D$ be a strict digraph — a loopless directed graph — on a finite, nonempty vertex
set. Write $\delta^-(D)$ for the minimum indegree and $\delta^+(D)$ for the minimum
outdegree, each taken over all vertices of $D$. Then $D$ contains a directed path — a
directed walk in which no vertex is repeated — whose length (its number of arcs) is at
least $\max\{\delta^-(D),\,\delta^+(D)\}$.
\end{theorem}

\begin{theorem}[A long directed cycle from large minimum degree]
\lean{exists_directedCycle_length_ge}
\label{exists_directedCycle_length_ge}
\leanok
\uses{Digraph.IsDirectedCycle, Digraph.IsStrict, Digraph.minIndegree, Digraph.minOutdegree, Digraph.toQuiver}
Let $D$ be a strict digraph on a finite, nonempty vertex set, and let $k$ be a positive
integer such that $\max\{\delta^{-}(D),\delta^{+}(D)\}=k$, where $\delta^{-}(D)$ and
$\delta^{+}(D)$ denote the minimum indegree and minimum outdegree of $D$. Then $D$
contains a directed cycle of length at least $k+1$.
\end{theorem}

\begin{theorem}[Powers of the adjacency matrix count directed walks]
\lean{adjMatrix_pow_apply_eq_card_directedWalk}
\label{adjMatrix_pow_apply_eq_card_directedWalk}
\leanok
\uses{Digraph.adjMatrix, Digraph.toQuiver}
Let $D$ be a digraph on a finite vertex set $V$, and let $A$ be its $0/1$ adjacency
matrix over $\bbn$, so $A_{uv} = 1$ when there is an arc from $u$ to $v$ and $A_{uv} =
0$ otherwise. Then for every $k \in \bbn$ and all vertices $u, v \in V$, the entry
$(A^k)_{uv}$ equals the number of directed walks of length $k$ from $u$ to $v$ in $D$,
that is, the number of sequences of $k$ arrows $u = w_0 \to w_1 \to \cdots \to w_k = v$.
\end{theorem}

\begin{theorem}[The condensation is acyclic]
\lean{condensation_acyclic}
\label{condensation_acyclic}
\leanok
\uses{Digraph.IsDirectedCycle, Digraph.condensation, Digraph.toQuiver}
Let $D$ be a digraph on a finite vertex set, and let $\hat{D}$ denote its condensation:
the digraph whose vertices are the strongly connected components (dicomponents) of $D$,
with an arc from one component to another precisely when $D$ has an arc whose tail lies
in the first and whose head lies in the second. Then $\hat{D}$ contains no directed
cycle; that is, for every dicomponent $c$ and every closed walk in $\hat{D}$ from $c$ to
$c$, that walk is not a directed cycle — it does not have positive length together with
distinct vertices apart from its coinciding endpoints.
\end{theorem}

\begin{theorem}[Existence of a balanced orientation]
\lean{exists_balanced_orientation}
\label{exists_balanced_orientation}
\leanok
\uses{Digraph.IsOrientationOf, Digraph.indegree, Digraph.outdegree}
Let $G$ be a finite simple graph. Then $G$ admits an orientation $D$ in which the in-
and outdegrees at every vertex differ by at most one; that is, $\abs{d^{+}(v) -
d^{-}(v)} \le 1$ for every vertex $v$, where $d^{+}(v)$ and $d^{-}(v)$ denote the
outdegree and indegree of $v$ in $D$.
\end{theorem}

\begin{theorem}[Gallai–Roy theorem: chromatic number forces a long directed path]
\lean{roy_gallai_directed_path}
\label{roy_gallai_directed_path}
\leanok
\uses{Digraph.IsDirectedPath, Digraph.toQuiver}
(Roy, 1967; Gallai, 1968.) Let $D$ be a digraph on a finite, nonempty vertex set $V$,
and suppose the underlying simple graph of $D$ has chromatic number $k$. Then $D$
contains a directed path — a directed walk $v_0, v_1, \dots, v_\ell$ whose consecutive
vertices are joined by arcs $v_{i-1} \to v_i$ and whose vertices are pairwise distinct —
of length at least $k - 1$, that is, with $\ell \ge k - 1$ arcs.
\end{theorem}

\begin{theorem}[Orientation with short directed paths for a k-colourable graph]
\lean{exists_orientation_longest_directedPath_le}
\label{exists_orientation_longest_directedPath_le}
\leanok
\uses{Digraph.IsDirectedPath, Digraph.IsOrientationOf, Digraph.toQuiver}
Let $G$ be a simple graph on a finite vertex set that admits a proper colouring with $k$
colours. Then $G$ has an orientation $D$ — a digraph whose underlying graph is $G$ and
in which no pair of vertices is joined by arcs in both directions — such that every
directed path in $D$ (a directed walk that repeats no vertex) has length at most $k-1$,
that is, passes through at most $k$ vertices.
\end{theorem}

\begin{theorem}[Rédei's theorem: every tournament has a directed Hamilton path]
\lean{redei_directed_hamilton_path}
\label{redei_directed_hamilton_path}
\leanok
\uses{Digraph.IsDirectedPath, Digraph.IsTournament, Digraph.toQuiver}
Let $D$ be a tournament on a finite, nonempty vertex set $V$ — a directed graph with no
loops in which, for every pair of distinct vertices $u$ and $v$, exactly one of the arcs
$u \to v$ and $v \to u$ is present. Then $D$ has a directed Hamilton path: there exist
vertices $u, v$ and a directed path from $u$ to $v$ that visits every vertex of $V$
exactly once.
\end{theorem}

\begin{theorem}[Existence of a quasi-kernel (Chvátal–Lovász)]
\lean{chvatal_lovasz_semikernel}
\label{chvatal_lovasz_semikernel}
\leanok
(Chvátal and Lovász, 1974.) Let $D$ be a loopless digraph on a finite vertex set $V$;
that is, no vertex has an arc to itself. Then $D$ admits an independent set $S \subseteq
V$ — no two distinct vertices of $S$ are joined by an arc of $D$ in either direction —
such that every vertex $v \notin S$ is reachable from some vertex $u \in S$ by a
directed path of length one or two: either $u \to v$ is an arc of $D$, or there is a
vertex $w$ with arcs $u \to w$ and $w \to v$.
\end{theorem}

\begin{theorem}[Every finite tournament has a king]
\lean{tournament_exists_king}
\label{tournament_exists_king}
\leanok
\uses{Digraph.IsTournament}
Let $D$ be a tournament on a finite, nonempty vertex set $V$; that is, for each pair of
distinct vertices exactly one of the two directed edges between them is present, and no
vertex has an edge to itself. Write $u \to v$ when the edge from $u$ to $v$ is present.
Then $D$ has a \emph{king}: a vertex $u$ such that every other vertex $v$ is reached
from $u$ by a directed path of length at most two — that is, either $u \to v$, or there
is a vertex $w$ with both $u \to w$ and $w \to v$.
\end{theorem}

\begin{theorem}[Reversing one arc to make a tournament diconnected]
\lean{tournament_diconnected_or_reorient_one}
\label{tournament_diconnected_or_reorient_one}
\leanok
\uses{Digraph.Diconnected, Digraph.IsTournament}
Let $D$ be a tournament on a finite, nonempty vertex set $V$. Then either $D$ is
diconnected, or there is an arc $x \to y$ of $D$ whose reversal makes it so: deleting
the arc $x \to y$ and inserting the arc $y \to x$, while leaving every other arc
unchanged, yields a diconnected digraph.
\end{theorem}

\begin{theorem}[Unilateral digraphs and spanning directed walks]
\lean{isUnilateral_iff_exists_spanning_directedWalk}
\label{isUnilateral_iff_exists_spanning_directedWalk}
\leanok
\uses{Digraph.IsUnilateral, Digraph.toQuiver}
Let $D$ be a digraph on a finite, nonempty vertex set $V$. Then $D$ is unilateral if and
only if there exist vertices $u, v \in V$ and a directed walk from $u$ to $v$ — a finite
sequence of edges, each traversed in the direction it points — whose vertices include
every vertex of $V$. In other words, $D$ is unilateral exactly when it admits a spanning
directed walk.
\end{theorem}

\begin{theorem}[Inserting a vertex into a non-extendable directed path in a tournament]
\lean{tournament_maximal_directedPath_insert}
\label{tournament_maximal_directedPath_insert}
\leanok
\uses{Digraph.IsTournament}
Let $D$ be a tournament on a finite vertex set $V$ — an orientation of the complete
graph on $V$. Let $v_1, v_2, \dots, v_k$ be a directed path in $D$: a nonempty list of
distinct vertices such that $v_j \to v_{j+1}$ is an arc of $D$ for every $1 \le j < k$.
Suppose this path cannot be extended at either end, meaning that for every vertex $w$
not among $v_1, \dots, v_k$ there is no arc from $w$ to $v_1$ and no arc from $v_k$ to
$w$. Then for every vertex $v$ not on the path there is an index $i$ with $1 \le i < k$
such that $v_i \to v$ and $v \to v_{i+1}$ are both arcs of $D$; that is, $v$ can be
inserted between two consecutive vertices of the path.
\end{theorem}

\begin{theorem}[Chvátal–Komlós monotone directed path theorem]
\lean{chvatal_komlos_monotone_directedPath}
\label{chvatal_komlos_monotone_directedPath}
\leanok
Let $D$ be a digraph on a finite vertex set $V$, let $f : V \to \bbr$ be a real-valued
function on the vertices, and let $m$ and $n$ be natural numbers. Write $\chi(D)$ for
the chromatic number of the underlying simple graph of $D$ (in which two vertices are
joined whenever there is an arc between them in either direction). If $\chi(D) > mn$,
then at least one of the following holds: either $D$ contains a directed path $u_0 \to
u_1 \to \cdots \to u_m$ on $m+1$ distinct vertices along which $f$ is nondecreasing,
$f(u_0) \le f(u_1) \le \cdots \le f(u_m)$; or $D$ contains a directed path $v_0 \to v_1
\to \cdots \to v_n$ on $n+1$ distinct vertices along which $f$ is strictly decreasing,
$f(v_0) > f(v_1) > \cdots > f(v_n)$.
\end{theorem}

\begin{theorem}[The Erdős–Szekeres theorem on monotone subsequences]
\lean{erdos_szekeres_of_chvatal_komlos}
\label{erdos_szekeres_of_chvatal_komlos}
\leanok
Let $m, n, N$ be natural numbers with $mn < N$, and let $g_0, g_1, \dots, g_{N-1}$ be a
sequence of $N$ distinct integers. Then the sequence has an increasing subsequence of
length $m$ — indices $i_1 < i_2 < \dots < i_m$ with $g_{i_1} < g_{i_2} < \dots <
g_{i_m}$ — or a decreasing subsequence of length $n$ — indices $j_1 < j_2 < \dots < j_n$
with $g_{j_1} > g_{j_2} > \dots > g_{j_n}$.
\end{theorem}

\begin{theorem}[Orientation with directed paths bounded by maximum degree]
\lean{exists_orientation_directedPath_le_maxDegree}
\label{exists_orientation_directedPath_le_maxDegree}
\leanok
\uses{Digraph.IsDirectedPath, Digraph.IsOrientationOf, Digraph.toQuiver}
Let $G$ be a finite simple graph with maximum degree $\Delta(G)$. Then $G$ admits an
orientation $D$ — a digraph on the same vertex set whose underlying graph is $G$ and in
which no two vertices are joined by arcs in both directions — such that every directed
path in $D$, that is, every directed walk that repeats no vertex, has length at most
$\Delta(G)$, where the length of a path is the number of arcs it traverses.
\end{theorem}

\begin{theorem}[Moon's theorem on vertex-pancyclicity of diconnected tournaments]
\lean{moon_vertex_pancyclic}
\label{moon_vertex_pancyclic}
\leanok
\uses{Digraph.Diconnected, Digraph.IsDirectedCycle, Digraph.IsTournament, Digraph.toQuiver}
Let $D$ be a diconnected tournament on a finite vertex set with at least three vertices,
and let $u$ be any vertex of $D$. Then for every integer $k$ with $3 \le k \le
\abs{V(D)}$ there is a directed cycle of length $k$ passing through $u$.
\end{theorem}

\begin{theorem}[Ghouila–Houri theorem for directed Hamilton cycles]
\lean{ghouila_houri_directed_hamilton_cycle}
\label{ghouila_houri_directed_hamilton_cycle}
\leanok
\uses{Digraph.IsDirectedCycle, Digraph.IsStrict, Digraph.indegree, Digraph.outdegree, Digraph.toQuiver}
Let $D$ be a loopless finite digraph on a vertex set of size $n > 2$, and suppose that
every vertex $v$ satisfies $d^{-}(v) \ge n/2$ and $d^{+}(v) \ge n/2$, where $d^{-}(v)$
and $d^{+}(v)$ denote the indegree and outdegree of $v$. Then $D$ contains a directed
cycle of length $n$; that is, a closed directed walk of positive length whose vertices,
apart from the coinciding endpoints, are distinct, and which passes through every vertex
of $D$.
\end{theorem}

\begin{theorem}[Existence of a directed Euler tour]
\lean{exists_directedEulerTour_iff}
\label{exists_directedEulerTour_iff}
\leanok
\uses{Digraph.IsDirectedEulerTour, Digraph.indegree, Digraph.outdegree, Digraph.toQuiver}
Let $D$ be a digraph on a finite, nonempty vertex set $V$ whose underlying undirected
graph is connected. Then $D$ admits a directed Euler tour—a closed directed trail that
traverses every arc of $D$ exactly once—if and only if every vertex $v \in V$ is
balanced, meaning its outdegree and indegree agree, $d^{+}(v) = d^{-}(v)$.
\end{theorem}

\begin{theorem}[Arc-disjoint directed paths from an excess to a deficit vertex]
\lean{exists_arcDisjoint_directedPaths}
\label{exists_arcDisjoint_directedPaths}
\uses{Digraph.IsDirectedPath, Digraph.arcsOf, Digraph.indegree, Digraph.outdegree, Digraph.toQuiver}
Let $D$ be a finite digraph, let $x \neq y$ be two of its vertices, and let $l$ be a
natural number. Suppose that $d^{+}(x) - d^{-}(x) = l$ and $d^{-}(y) - d^{+}(y) = l$,
and that every other vertex $v$ (that is, $v \neq x$ and $v \neq y$) satisfies $d^{+}(v)
= d^{-}(v)$, where $d^{+}$ and $d^{-}$ denote outdegree and indegree. Then $D$ contains
$l$ pairwise arc-disjoint directed paths from $x$ to $y$.
\end{theorem}

\begin{theorem}[Odd cycle forces an odd directed cycle in a diconnected digraph]
\lean{exists_directed_odd_cycle}
\label{exists_directed_odd_cycle}
\leanok
\uses{Digraph.Diconnected, Digraph.IsDirectedCycle, Digraph.toQuiver}
Let $D$ be a diconnected digraph on a finite vertex set, so that every vertex is
reachable from every other along directed paths. Suppose the underlying simple graph of
$D$, obtained by forgetting the orientation of each arc, contains a cycle of odd length.
Then $D$ contains a directed cycle of odd length.
\end{theorem}

\begin{theorem}[Diconnection equals 1-arc-connectedness]
\lean{diconnected_iff_isKArcConnected_one}
\label{diconnected_iff_isKArcConnected_one}
\leanok
\uses{Digraph.Diconnected, Digraph.IsKArcConnected}
Let $D$ be a finite digraph whose vertex set has at least two vertices. Then $D$ is
diconnected — every vertex is reachable from every other along a directed path — if and
only if $D$ is $1$-arc-connected; that is, if and only if for every nonempty proper
subset $S$ of the vertex set there is at least one arc of $D$ with tail in $S$ and head
outside $S$.
\end{theorem}

\begin{theorem}[Arc-connectivity of the associated digraph equals edge connectivity]
\lean{associatedDigraph_isKArcConnected_iff}
\label{associatedDigraph_isKArcConnected_iff}
\leanok
\uses{Digraph.IsKArcConnected, SimpleGraph.associatedDigraph, edgeConnectivity}
Let $G$ be a finite simple graph on a vertex set with at least two vertices, let $D(G)$
be its associated digraph (each edge of $G$ replaced by the two oppositely oriented arcs
joining its ends), and let $k$ be a natural number. Then $D(G)$ is $k$-arc-connected if
and only if $k \le \kappa'(G)$, where $\kappa'(G)$ denotes the edge connectivity of $G$.
\end{theorem}

\begin{theorem}[Every vertex of the de Bruijn digraph has indegree and outdegree two]
\lean{deBruijnDigraph_indegree_eq_two}
\label{deBruijnDigraph_indegree_eq_two}
\leanok
\uses{Digraph.indegree, Digraph.outdegree, deBruijnDigraph}
Let $n \ge 2$ be an integer, and let $v$ be any vertex of the de Bruijn digraph $D_n$.
Then both the indegree and the outdegree of $v$ equal $2$; that is, $d^{-}(v) = 2$ and
$d^{+}(v) = 2$.
\end{theorem}

\begin{theorem}[Connectivity of the de Bruijn digraph]
\lean{deBruijnDigraph_connected}
\label{deBruijnDigraph_connected}
\leanok
\uses{deBruijnDigraph}
Fix an integer $n \ge 2$, and let $D_n$ be the de Bruijn digraph whose vertices are the
binary strings of length $n-1$ and whose arcs run from each string $p_1 p_2 \cdots
p_{n-1}$ to every string of the form $p_2 p_3 \cdots p_{n-1} b$ with $b \in \{0,1\}$.
Form the underlying undirected simple graph on the same vertex set, in which two
distinct strings are adjacent precisely when an arc of $D_n$ joins one to the other in
either direction. Then this graph is connected: its vertex set is nonempty and any two
vertices are linked by a path.
\end{theorem}

\begin{theorem}[Existence of a directed Euler tour in the de Bruijn digraph]
\lean{deBruijnDigraph_exists_directedEulerTour}
\label{deBruijnDigraph_exists_directedEulerTour}
\uses{Digraph.IsDirectedEulerTour, Digraph.toQuiver, deBruijnDigraph}
For every integer $n \ge 2$, the de Bruijn digraph $D_n$ admits a directed Euler tour:
there exist a vertex $u$ and a closed directed walk based at $u$ that traverses each arc
of $D_n$ exactly once.
\end{theorem}

\begin{theorem}[Robbins' theorem on diconnected orientations]
\lean{robbins_orientation}
\label{robbins_orientation}
\leanok
\uses{Digraph.Diconnected, Digraph.IsOrientationOf, edgeConnectivity}
Let $G$ be a finite simple graph on a nonempty vertex set, and suppose its edge
connectivity satisfies $\kappa'(G) \ge 2$. Then $G$ has a diconnected orientation: one
can assign a direction to each edge of $G$, forming a digraph $D$ whose underlying graph
is $G$, so that for every ordered pair of vertices $u, v$ there is a directed path in
$D$ from $u$ to $v$.
\end{theorem}

\begin{theorem}[A $k$-arc-connected orientation of a $2k$-edge-connected Eulerian graph]
\lean{exists_kArcConnected_orientation_of_eulerian}
\label{exists_kArcConnected_orientation_of_eulerian}
\leanok
\uses{Digraph.IsKArcConnected, Digraph.IsOrientationOf, edgeConnectivity}
Let $G$ be a finite graph on a nonempty vertex set that admits an Eulerian trail — a
walk traversing every edge of $G$ exactly once — and let $k$ be a natural number with
$2k \le \kappa'(G)$. Then $G$ has an orientation that is $k$-arc-connected.
\end{theorem}

\begin{theorem}[Positivity of a diameter-shifted power of a tournament's adjacency matrix]
\lean{tournament_adjMatrix_pow_pos}
\label{tournament_adjMatrix_pow_pos}
\leanok
\uses{Digraph.Diconnected, Digraph.IsTournament, Digraph.adjMatrix, Digraph.dirDiameter}
Let $D$ be a diconnected tournament on a finite vertex set $V$ with $\nu = \abs{V} \ge
5$, and let $A$ be its $\nu \times \nu$ adjacency matrix over $\bbn$, so that $A_{uv} =
1$ when there is an arc from $u$ to $v$ and $A_{uv} = 0$ otherwise. Writing $d$ for the
directed diameter of $D$, every entry of the matrix power $A^{\,d+3}$ is strictly
positive; that is, $\bigl(A^{\,d+3}\bigr)_{ij} > 0$ for all $i, j \in V$.
\end{theorem}

\begin{theorem}[Primitivity of a tournament's adjacency matrix]
\lean{tournament_adjMatrix_isPrimitive_iff}
\label{tournament_adjMatrix_isPrimitive_iff}
\leanok
\uses{Digraph.Diconnected, Digraph.IsTournament, Digraph.adjMatrix, Matrix.IsPrimitive}
Let $D$ be a tournament on a finite vertex set $V$ (an orientation of the complete graph
on $V$), and let $A$ be its $0/1$ adjacency matrix. Then $A$ is primitive — that is,
some power $A^{k}$ has all entries positive — if and only if $D$ is diconnected (every
vertex is reachable from every other) and $V$ has at least four vertices.
\end{theorem}
